% !TEX root = ../main.tex
% Figure: Epitaxial stack cross-section with layer functions
% Chapter: ch16 - Epitaxial and Optical Structure Design
% Description: Complete vertical structure of a typical PCSEL

\begin{figure}[htbp]
  \centering
  \begin{tikzpicture}[scale=1.0]
    % Substrate at bottom
    \fill[gray!40] (-3,-5) rectangle (3,-4);
    \draw[thick] (-3,-4) -- (3,-4);
    \node[right] at (3.1,-4.5) {GaAs 衬底 ($n$-doped)};
    \node[left, gray!60] at (-3,-4.5) {\scriptsize $500\,\mu\mathrm{m}$};
    
    % Lower DBR mirror
    \foreach \i in {0,...,7} {
      \pgfmathparse{\i*0.3-3.8}
      \fill[blue!\the\numexpr\i*10!green!60!black!20] (-3,\pgfmathresult) rectangle (3,{\pgfmathresult+0.15});
      \draw[thin, gray] (-3,\pgfmathresult) -- (3,\pgfmathresult);
    }
    \draw[thick] (-3,-3.8) -- (3,-3.8);
    \node[right] at (3.1,-3.5) {$n$-DBR (20--30 对)};
    \node[left, blue] at (-3,-3.5) {\scriptsize Al$_{0.9}$Ga$_{0.1}$As/Al$_{0.2}$Ga$_{0.8}$As};
    
    % Lower SCH waveguide
    \fill[green!15] (-3,-3.2) rectangle (3,-2.5);
    \draw[thick] (-3,-3.2) -- (3,-3.2);
    \draw[thick] (-3,-2.5) -- (3,-2.5);
    \node[right] at (3.1,-2.85) {下 SCH 波导层};
    \node[left, green!60!black] at (-3,-2.85) {\scriptsize $200\,\mathrm{nm}$};
    
    % Active region (MQW)
    \fill[red!20] (-3,-2.5) rectangle (3,-2.2);
    \draw[thick] (-3,-2.5) -- (3,-2.5);
    % Quantum wells inside active region
    \foreach \y in {-2.45,-2.35,-2.25} {
      \fill[red!60] (-3,\y) rectangle (3,{\y+0.05});
      \draw[dashed, red] (-3,\y) -- (3,\y);
    }
    \draw[thick] (-3,-2.2) -- (3,-2.2);
    \node[right] at (3.1,-2.35) {有源区 (3--5 阱 MQW)};
    \node[left, red] at (-3,-2.35) {\scriptsize In$_{0.2}$Ga$_{0.8}$As};
    
    % Upper SCH waveguide
    \fill[green!15] (-3,-2.2) rectangle (3,-1.5);
    \draw[thick] (-3,-2.2) -- (3,-2.2);
    \draw[thick] (-3,-1.5) -- (3,-1.5);
    \node[right] at (3.1,-1.85) {上 SCH 波导层};
    \node[left, green!60!black] at (-3,-1.85) {\scriptsize $200\,\mathrm{nm}$};
    
    % Photonic crystal layer
    \fill[yellow!30] (-3,-1.5) rectangle (3,-0.8);
    \draw[thick] (-3,-1.5) -- (3,-1.5);
    % Etched holes
    \foreach \x in {-2.5,-1.5,...,2.5} {
      \fill[white] (\x,-1.15) circle (0.15);
      \draw[thick] (\x,-1.15) circle (0.15);
    }
    \draw[thick] (-3,-0.8) -- (3,-0.8);
    \node[right] at (3.1,-1.15) {PhC 光栅层 (刻蚀深度~$150\,\mathrm{nm}$)};
    \node[left, yellow!60!black] at (-3,-1.15) {\scriptsize $\Lambda=260\,\mathrm{nm}$};
    
    % Upper cladding / contact
    \fill[orange!20] (-3,-0.8) rectangle (3,-0.3);
    \draw[thick] (-3,-0.8) -- (3,-0.8);
    \draw[thick] (-3,-0.3) -- (3,-0.3);
    \node[right] at (3.1,-0.55) {$p$-AlGaAs 包层};
    \node[left, orange] at (-3,-0.55) {\scriptsize $500\,\mathrm{nm}$};
    
    % Contact metal
    \fill[gray!60] (-3,-0.3) rectangle (3,0);
    \draw[thick] (-3,-0.3) -- (3,-0.3);
    \draw[thick] (-3,0) -- (3,0);
    \node[right] at (3.1,-0.15) {$p$-接触电极 (Ti/Pt/Au)};
    \node[left, gray] at (-3,-0.15) {\scriptsize $200\,\mathrm{nm}$};
    
    % Bonding layer (optional)
    \fill[purple!20] (-3,0) rectangle (3,0.3);
    \draw[thick] (-3,0) -- (3,0);
    \draw[thick] (-3,0.3) -- (3,0.3);
    \node[right] at (3.1,0.15) {键合层 (Au-Sn/$\mathrm{SiO_2}$)};
    \node[left, purple] at (-3,0.15) {\scriptsize 倒装焊可选};
    
    % Heat sink at top
    \fill[cyan!30] (-3,0.3) rectangle (3,1);
    \draw[thick] (-3,0.3) -- (3,0.3);
    \draw[thick] (-3,1) -- (3,1);
    \node[right] at (3.1,0.65) {热沉 (Cu/W 或金刚石)};
    \node[left, cyan!60!black] at (-3,0.65) {\scriptsize 高热导率};
    
    % Current flow arrows
    \draw[->, very thick, red] (-4,-0.15) -- node[left] {$I_{inj}$} (-4,-3.5);
    
    % Light emission arrow
    \draw[->, very thick, green!60!black] (0,-0.8) -- (0,1.5);
    \node[right, green!60!black, font=\bfseries] at (0.2,1.2) {激光输出};
    
    % Total thickness bracket
    \draw[thick] (3.5,-4) -- node[right] {外延总厚~$8$--$12\,\mu\mathrm{m}$} (3.5,0.3);
    
    % Mode profile overlay (schematic)
    \draw[very thick, blue, opacity=0.5, domain=-3:3, samples=100] 
      plot(\x,{0.3*cos(deg(0.5*\x))*exp(-(\x+0.5)*(\x+0.5)/4)-1.85});
    \node[blue, align=center] at (2,-1.5) {基模场分布\\$\mathrm{LP}_{01}$};
  \end{tikzpicture}
  \caption{典型 PCSEL 外延堆栈的截面示意图（未按比例）。图中依次给出衬底、DBR、上下 SCH 波导层、有源量子阱区、光子晶体层、包层与金属接触，并标出注入电流方向、垂直出光方向以及基模近场包络。右侧括号给出总外延厚度的典型量级。}
  \label{fig:ch16_epitaxial_stack}
\end{figure}
